\section{Comparison by terrain class}\label{sec:strat}

Does the difference come from the ground the two disciplines fly over?

Hang gliders fly higher ground than paragliders (Fig.~\ref{fig:hist}), and $H$ depends on the ground. The comparison is
therefore made within each class of Section~\ref{sec:terrain}, on the flights whose class is known. The classes with fewer
than \UniMinCell{} hang glider flights would not be fitted; none is in this case.

\paragraph{A, fixed population.} At $T=$ \UniTauFixed{} s (Fig.~\ref{fig:strat-A}, Table~\ref{tab:strat-A}) the paragliders have
the larger $H$ in every class: $\Hpara-\Hhang=$ \UniDHPlainsAfixed{} \UniDHBandPlainsAfixed{} in plains,
\UniDHHillsAfixed{} \UniDHBandHillsAfixed{} in hills, \UniDHLowAfixed{} \UniDHBandLowAfixed{} in low mountains and
\UniDHHighAfixed{} \UniDHBandHighAfixed{} in high mountains. $H$ falls with the ruggedness of the ground in both disciplines: from
\UniHFirstPlainsAfixed{} in plains to \UniHFirstHighAfixed{} in high mountains for paragliders, from \UniHSecondPlainsAfixed{} to
\UniHSecondHighAfixed{} for hang gliders. The raw difference (\UniDHAllAfixed) lies within the range of the differences by class, so the terrain mix does not
account for it.

\begin{figure}
\centering
\includegraphics[width=\textwidth]{strat_A_fixed.pdf}
\caption{Within every class the curve of the paragliders lies above that of the hang gliders at lags up to 1000 s, the two
curves overlapping most in plains. Design A, $T=$ \UniTauFixed{} s. Panels as Fig.~\ref{fig:raw-A}; flights above each panel.}\label{fig:strat-A}
\end{figure}

With the longest $T$ of each class (Fig.~\ref{fig:strat-A-longest}) the differences are \UniDHPlainsAlongest{} (plains, $T=$
\UniTPlainsAlongest{} s), \UniDHHillsAlongest{} (hills, $T=$ \UniTHillsAlongest{} s, the same as the fixed one), \UniDHLowAlongest{} (low mountains,
$T=$ \UniTLowAlongest{} s) and \UniDHHighAlongest{} (high mountains, $T=$ \UniTHighAlongest{} s). In the two mountain classes
$\alpha(\tau)$ falls above 1000 s for both disciplines and the hang curve ends lower.

\begin{figure}
\centering
\includegraphics[width=\textwidth]{strat_A_longest.pdf}
\caption{With the longest $T$ of each class the difference persists in the mountains, and the bands of the hang gliders widen
above 1000 s. Design A, $T$ in the panel labels.}\label{fig:strat-A-longest}
\end{figure}

\begin{table}
\centering
\small
\renewcommand{\arraystretch}{1.15}\setlength{\tabcolsep}{4pt}
\begin{adjustbox}{max width=\textwidth}
\begin{tabular}{lrrccc}
\toprule
Population & flights (pilots) para & flights (pilots) hang & $H$ para & $H$ hang & $\Delta H$ \\
\midrule
plains, A, $T=1000$ s & 11\,990 (1\,368) & 170 (34) & \shortstack{0.941\\[-1pt]{\scriptsize [0.940, 0.942]}} & \shortstack{0.932\\[-1pt]{\scriptsize [0.923, 0.941]}} & \shortstack{$+$0.009\\[-1pt]{\scriptsize [$+$0.000, $+$0.018]}} \\
plains, A, longest $T$ (4150 s) & 6\,541 (981) & 117 (29) & \shortstack{0.946\\[-1pt]{\scriptsize [0.944, 0.947]}} & \shortstack{0.933\\[-1pt]{\scriptsize [0.925, 0.941]}} & \shortstack{$+$0.013\\[-1pt]{\scriptsize [$+$0.005, $+$0.021]}} \\
\midrule
hills, A, $T=1000$ s & 10\,218 (1\,900) & 100 (42) & \shortstack{0.927\\[-1pt]{\scriptsize [0.926, 0.928]}} & \shortstack{0.909\\[-1pt]{\scriptsize [0.902, 0.916]}} & \shortstack{$+$0.018\\[-1pt]{\scriptsize [$+$0.010, $+$0.025]}} \\
hills, A, longest $T$ (1000 s) & 10\,218 (1\,900) & 100 (42) & \shortstack{0.927\\[-1pt]{\scriptsize [0.926, 0.928]}} & \shortstack{0.909\\[-1pt]{\scriptsize [0.902, 0.916]}} & \shortstack{$+$0.018\\[-1pt]{\scriptsize [$+$0.010, $+$0.025]}} \\
\midrule
low mountains, A, $T=1000$ s & 20\,841 (4\,064) & 775 (150) & \shortstack{0.909\\[-1pt]{\scriptsize [0.908, 0.910]}} & \shortstack{0.882\\[-1pt]{\scriptsize [0.877, 0.887]}} & \shortstack{$+$0.027\\[-1pt]{\scriptsize [$+$0.022, $+$0.032]}} \\
low mountains, A, longest $T$ (9350 s) & 5\,188 (1\,871) & 113 (57) & \shortstack{0.902\\[-1pt]{\scriptsize [0.900, 0.904]}} & \shortstack{0.861\\[-1pt]{\scriptsize [0.852, 0.871]}} & \shortstack{$+$0.040\\[-1pt]{\scriptsize [$+$0.030, $+$0.050]}} \\
\midrule
high mountains, A, $T=1000$ s & 8\,774 (2\,691) & 224 (81) & \shortstack{0.907\\[-1pt]{\scriptsize [0.906, 0.908]}} & \shortstack{0.880\\[-1pt]{\scriptsize [0.873, 0.887]}} & \shortstack{$+$0.026\\[-1pt]{\scriptsize [$+$0.019, $+$0.034]}} \\
high mountains, A, longest $T$ (7630 s) & 4\,026 (1\,597) & 109 (47) & \shortstack{0.897\\[-1pt]{\scriptsize [0.895, 0.899]}} & \shortstack{0.869\\[-1pt]{\scriptsize [0.857, 0.881]}} & \shortstack{$+$0.028\\[-1pt]{\scriptsize [$+$0.016, $+$0.040]}} \\
\bottomrule
\end{tabular}
\end{adjustbox}
\caption{As Table~\ref{tab:raw}, design A per terrain class.}\label{tab:strat-A}
\end{table}


\paragraph{B, all flights.} The same comparison on all flights (Fig.~\ref{fig:strat-B}, Table~\ref{tab:strat-B}) gives, up to 1000 s,
\UniDHPlainsBkilo, \UniDHHillsBkilo, \UniDHLowBkilo{} and \UniDHHighBkilo{} in the four classes and, up to the longest $T$ of the class,
\UniDHPlainsBlongest, \UniDHHillsBlongest, \UniDHLowBlongest{} and \UniDHHighBlongest.

\begin{figure}
\centering
\includegraphics[width=\textwidth]{strat_B.pdf}
\caption{In every class the number of flights falls at long lags and the curves bend away from those of the fixed population there.
Design B; panels as Fig.~\ref{fig:raw-B}.}\label{fig:strat-B}
\end{figure}

\begin{table}
\centering
\small
\renewcommand{\arraystretch}{1.15}\setlength{\tabcolsep}{4pt}
\begin{adjustbox}{max width=\textwidth}
\begin{tabular}{lrrccc}
\toprule
Population & flights (pilots) para & flights (pilots) hang & $H$ para & $H$ hang & $\Delta H$ \\
\midrule
plains, B, to 1000 s & 11\,990 (1\,368) & 170 (34) & \shortstack{0.941\\[-1pt]{\scriptsize [0.940, 0.942]}} & \shortstack{0.931\\[-1pt]{\scriptsize [0.921, 0.940]}} & \shortstack{$+$0.010\\[-1pt]{\scriptsize [$+$0.000, $+$0.020]}} \\
plains, B, to longest $T$ (4150 s) & 6\,541 (981) & 117 (29) & \shortstack{0.944\\[-1pt]{\scriptsize [0.942, 0.945]}} & \shortstack{0.931\\[-1pt]{\scriptsize [0.923, 0.939]}} & \shortstack{$+$0.013\\[-1pt]{\scriptsize [$+$0.005, $+$0.021]}} \\
\midrule
hills, B, to 1000 s & 10\,218 (1\,900) & 100 (42) & \shortstack{0.927\\[-1pt]{\scriptsize [0.926, 0.928]}} & \shortstack{0.909\\[-1pt]{\scriptsize [0.902, 0.917]}} & \shortstack{$+$0.017\\[-1pt]{\scriptsize [$+$0.010, $+$0.025]}} \\
hills, B, to longest $T$ (1000 s) & 10\,218 (1\,900) & 100 (42) & \shortstack{0.927\\[-1pt]{\scriptsize [0.926, 0.928]}} & \shortstack{0.909\\[-1pt]{\scriptsize [0.902, 0.917]}} & \shortstack{$+$0.017\\[-1pt]{\scriptsize [$+$0.010, $+$0.025]}} \\
\midrule
low mountains, B, to 1000 s & 20\,841 (4\,064) & 775 (150) & \shortstack{0.909\\[-1pt]{\scriptsize [0.908, 0.910]}} & \shortstack{0.882\\[-1pt]{\scriptsize [0.877, 0.887]}} & \shortstack{$+$0.027\\[-1pt]{\scriptsize [$+$0.022, $+$0.032]}} \\
low mountains, B, to longest $T$ (9350 s) & 5\,188 (1\,871) & 113 (57) & \shortstack{0.893\\[-1pt]{\scriptsize [0.891, 0.895]}} & \shortstack{0.851\\[-1pt]{\scriptsize [0.844, 0.858]}} & \shortstack{$+$0.042\\[-1pt]{\scriptsize [$+$0.035, $+$0.049]}} \\
\midrule
high mountains, B, to 1000 s & 8\,774 (2\,691) & 224 (81) & \shortstack{0.907\\[-1pt]{\scriptsize [0.906, 0.908]}} & \shortstack{0.880\\[-1pt]{\scriptsize [0.873, 0.887]}} & \shortstack{$+$0.027\\[-1pt]{\scriptsize [$+$0.020, $+$0.034]}} \\
high mountains, B, to longest $T$ (7630 s) & 4\,026 (1\,597) & 109 (47) & \shortstack{0.889\\[-1pt]{\scriptsize [0.887, 0.891]}} & \shortstack{0.861\\[-1pt]{\scriptsize [0.849, 0.874]}} & \shortstack{$+$0.028\\[-1pt]{\scriptsize [$+$0.015, $+$0.041]}} \\
\bottomrule
\end{tabular}
\end{adjustbox}
\caption{As Table~\ref{tab:raw}, design B per terrain class.}\label{tab:strat-B}
\end{table}


\keep{In each class of ground the paragliders have the larger $H$ at 1000 s, $\Hpara-\Hhang$ going from \UniDHPlainsAfixed{} to \UniDHLowAfixed{}; the
terrain mix does not account for the raw difference. In both disciplines $H$ falls from plains to mountains.}
